\subsection{Discrete mixed-characteristic bases with finitely generated residue field}

\begin{claim}{B.META.20}{Proved}{Completion, Kummer kernels and the deep-unit subspace beyond finite local fields}
Let $D$ be finitely generated over $\mathbb Q$, with
a discrete valuation of mixed characteristic $(0,p)$
whose residue field $k$ is finitely generated over
$\mathbb F_p$.
Let $M$ be its completion, choose an algebraic closure
$\overline M$, and let $\Omega$ be the algebraic closure
of $D$ inside $\overline M$.
The following assertions hold, also after every finite
extension of the valued field $D$.

First, $D^h=M\cap\Omega$, and completion gives an
equivalence on finite separable extensions of $D^h$
and $M$.  Thus $G_M$ identifies with the chosen
decomposition subgroup of $G_D$, with its inertia
subgroup $I_M$.
For finite $L/M$, its algebraic part $L\cap\Omega$
is finite over $D^h$, dense in $L$, and $L$ has an
$M$-basis in that algebraic part.
The field $\Omega$ admits a complex embedding.

Second, the rational $p$-Kummer kernel of $M$ is
exactly the finite group $\mu(M)$.
There is a canonical injective $\mathbb Q_p$-linear map
\[
         \iota_M:M\lhook\joinrel\longrightarrow
                      H^1(G_M,\mathbb Q_p(1)),                 \tag{META.DV1}
\]
characterized on sufficiently small $x$ by
$\iota_M(x)=\kappa_M(\exp x)$.
It is compatible with finite extensions and their
Galois actions.
No assertion that (META.DV1) exhausts the cohomology
is made.

Finally the image of $I_M$ in $G_{\mathbb Q_p}$ is
$I_{L_0}$ for some finite extension $L_0/\mathbb Q_p$.
In particular the cyclotomic character on $I_M$ is
nontrivial, restriction
$H^1(G_M,\mathbb Q_p(1))\to H^1(I_M,\mathbb Q_p(1))$
is injective, and every successive closed derived
quotient of $I_M$ is nontrivial.
\end{claim}
\begin{proof}
Write $v_M(M^\times)=\mathbb Z$ and $e=v_M(p)>0$.
The completion is a complete discrete valuation
field with the same value group and residue field.
The real-valued henselian comparison in Temkin,
Theorem~1.9.5(ii), \S1.9.7 and Lemma~1.9.8,
identifies the henselization with the relative
separable closure in the completion
\citep{AN_TemkinValuations}.
In characteristic zero this gives $D^h=M\cap\Omega$.
The proof of Stacks Tag~09EJ, Lemma~15.115.2,
works for an arbitrary DVR and supplies descent
of every finite separable extension of $M$
\citep{AN_Stacks09EJ}.
Here is also the full-faithfulness comparison.
A henselian field is relatively separably closed
in its completion.  Since $D^h$ is relatively
algebraically closed in $M$, it is linearly disjoint
from $M$ in each finite Galois extension; hence
completion preserves finite extension degrees.
Every finite extension of $D^h$ is henselian,
and its field embeddings preserve the uniquely
extended valuation, so extend uniquely to embeddings
of its completion.
Conversely an embedding of those completions carries
the elements algebraic over $D^h$ into the relative
algebraic part of the target, which is exactly the
original finite henselian field.
This proves full faithfulness and the claimed
equivalence.  It identifies the Galois and inertia
groups by preserving the residue extension.
It also proves the assertion about $L\cap\Omega$;
any basis before completion remains a basis afterwards.
In particular the algebraic part is dense.
Finally an embedding of the algebraic closure of
$\mathbb Q$ in $\Omega$ into $\mathbb C$ extends
across a finite transcendence basis and then across
its algebraic closure, giving the asserted embedding.
These arguments used neither finiteness of the
residue field nor local compactness.

We next prove the Kummer statements, including the
distinction between deep units and all units.
The algebraic closure $k_0$ of $\mathbb F_p$ in $k$
is finite.  Indeed if $t$ is a finite transcendence
basis, the degrees of all finite constant subfields
are bounded by $[k:\mathbb F_p(t)]$.
Moreover
\[
                       \bigcap_{n\ge0}k^{p^n}=k_0.             \tag{META.DV2}
\]
To see this, take a normal proper model of $k/k_0$.
The divisor of an element in every $k^{p^n}$ has
all its integer coefficients divisible by every
$p^n$, and is therefore zero.
Normality makes that rational function regular
everywhere, and properness makes it algebraic over
$k_0$.  It consequently belongs to $k_0$.
The reverse inclusion holds because the finite
field $k_0$ is perfect.

Put $U_m=1+\mathfrak m_M^m$.
For an integer $m>e/(p-1)$ the convergent power
series for logarithm and exponential give inverse
group isomorphisms
\[
             U_m\xrightarrow{\log}\mathfrak m_M^m,
                 \qquad\log(u^p)=p\log u.                    \tag{META.DV3}
\]
For completeness, in the logarithm the terms
$x^j/j$ tend to zero for $v_M(x)>0$; for the
exponential, $v_M(j!)\le ej/(p-1)$ gives convergence
when $v_M(x)>e/(p-1)$.
The first term has strictly smaller valuation
than all subsequent terms on that smaller ball;
the formal inverse identities then converge and
give (META.DV3).
Repeated $p$-th powers carry all of $U_1$ into
$U_m$ after a uniformly bounded number $N$ of
iterations.  In fact, if $v_M(u-1)=a\ge1$, the
binomial expansion gives
\[
                  v_M(u^p-1)\ge\min(pa,a+e)>a.
\]
Thus $N=m$ is already a sufficient bound.

The topology on $U_m$ induced by the $p$-power
quotients of $M^\times$ is its full $p$-adic topology.
If $b^{p^n}\in U_m$, valuation gives $v_M(b)=0$,
and reduction gives $\overline b^{p^n}=1$.
In characteristic $p$ this forces $\overline b=1$,
so $b\in U_1$.
For $n\ge N$ we obtain the explicit inclusions
\[
 U_m^{p^n}\ \subset\ U_m\cap(M^\times)^{p^n}
                 \ \subset\ U_m^{p^{\,n-N}}.                 \tag{META.DV4}
\]
Completeness and (META.DV3) now give an injection
\[
 U_m=\varprojlim_n U_m/U_m^{p^n}
       \lhook\joinrel\longrightarrow
             \varprojlim_n M^\times/(M^\times)^{p^n}.
\]
The inverse-limit Kummer sequence identifies the
right side with $H^1(G_M,\mathbb Z_p(1))$.
Its possible $\varprojlim^1$ term is zero:
the groups of $p^n$-th roots of unity in $M$ are
finite, so their inverse system is Mittag--Leffler.
Continuous cocycles of a profinite group in
$\mathbb Q_p(1)$ have bounded image; multiplication
by a power of $p$ puts them in $\mathbb Z_p(1)$.
The same observation for coboundaries proves
\[
 H^1(G_M,\mathbb Z_p(1))\otimes_{\mathbb Z_p}\mathbb Q_p
                         =H^1(G_M,\mathbb Q_p(1)).
\]
After tensoring the injection with $\mathbb Q_p$,
(META.DV3) identifies its domain with $M$.
This is (META.DV1); explicitly, for any $x\in M$
and sufficiently large $r$,
\[
                \iota_M(x)=p^{-r}\kappa_M(\exp(p^rx)).
\]
The formula proves independence of $r$, as well
as extension compatibility and Galois equivariance.

Suppose $\kappa_M(a)=0$.
Some $p^j$ multiple of the integral Kummer class
is zero, so $a^{p^j}\in(M^\times)^{p^n}$ for every $n$.
The value group gives $v_M(a)=0$, and reduction
together with (META.DV2) puts $\overline a$ in
$k_0^\times$.
Hensel's lemma lifts it to a root of unity $\zeta$
of order prime to $p$.
Then $u=a/\zeta\in U_1$ and $\kappa_M(u)=0$.
Raise it to $p^N$ to enter $U_m$.
The injection just proved gives $\log(u^{p^N})=0$,
and (META.DV3) gives $u^{p^N}=1$.
Thus $a$ is a root of unity; the converse follows
by tensoring the torsion Kummer class with $\mathbb Q_p$.
There are only finitely many roots of unity in $M$:
the prime-to-$p$ ones inject into the finite group
$k_0^\times$, while a primitive $p^n$-th root has
normalized $p$-adic value
$v_p(\zeta_{p^n}-1)=1/(p^{n-1}(p-1))$.
The latter follows from the Eisenstein polynomial
$\Phi_{p^n}(1+T)$, and those denominators must divide
$e$.  This also includes $p=2$.
All these assertions apply to finite extensions,
whose residue fields remain finitely generated.

It remains to justify the useful finite-local
quotient of inertia without asserting wild freeness.
Embed $\overline{\mathbb Q}_p$ in $\overline M$
over the canonical copy $\mathbb Q_p\subset M$,
and put
\[
                  T=M^{\mathrm{nr}}\cap\overline{\mathbb Q}_p.
\]
Finite unramified constant extensions stay etale
after base change to the valuation ring of $M$,
so $\mathbb Q_p^{\mathrm{nr}}\subset T$.
Every finite $L'/\mathbb Q_p$ inside $T$ has
$e(L'/\mathbb Q_p)\mid e$, since $M^{\mathrm{nr}}$
has the same value group as $M$.
Also
\[
 [L'\mathbb Q_p^{\mathrm{nr}}:\mathbb Q_p^{\mathrm{nr}}]
                                      =e(L'/\mathbb Q_p).
\]
For this equality, take the maximal unramified
subextension of $L'/\mathbb Q_p$ and an Eisenstein
polynomial for a uniformizer of the remaining
totally ramified extension.
After unramified base change the polynomial remains
Eisenstein, giving exactly the displayed degree.
The union $T/\mathbb Q_p^{\mathrm{nr}}$ thus has
bounded finite subextension degree, at most $e$,
and is finite.  Hence $T=L_0\mathbb Q_p^{\mathrm{nr}}$
for one finite $L_0/\mathbb Q_p$.
The image of $I_M$ in $G_{\mathbb Q_p}$ is closed
and has fixed field $T$, so is exactly $I_{L_0}$.
The local cyclotomic character has infinite image
there.  Thus $\mathbb Q_p(1)^{I_M}=0$, and
inflation--restriction proves the asserted injection.

Finally $I_{L_0}$ has strictly decreasing closed
derived series.  Its wild inertia is free pro-$p$
of infinite rank by Iwasawa Theorem~2(ii)
\citep{AN_Iwasawa1955}.
For every $r$, mapping finitely many free generators
to the adjacent elementary matrices of
$\operatorname{UT}_{2^r+1}(\mathbb F_p)$ shows that
its $r$-th derived subgroup is nontrivial:
balanced commutators use
$[1+E_{ab},1+E_{bc}]=1+E_{ac}$ and terminate at
$1+E_{1,2^r+1}$.
Hence every derived subgroup of $I_{L_0}$ is
nontrivial.  From the first commutator onward it
is pro-$p$, and a nontrivial pro-$p$ group has
a nontrivial abelian quotient, by taking any
nontrivial finite $p$-group quotient and then
a quotient of order $p$.
The first derived quotient is already nontrivial
by tame inertia.
Surjectivity $I_M\to I_{L_0}$ carries closed
derived subgroups onto closed derived subgroups,
so proves strictness for $I_M$ as well.
\end{proof}

\begin{claim}{B.META.21}{Proved}{Reconstruction at derived and nilpotent inertia fields with arbitrary finite residue transcendence}
Keep $D,M,\Omega,k$ of Claim~B.META.20, and put
$I=I_M\subset G_D$.
For the closed derived series $I^{(0)}=I$ and
$I^{(r+1)}=\overline{[I^{(r)},I^{(r)}]}$, and the
closed lower central series $\gamma_c(I)$, set
\[
 F_r=\Omega^{I^{(r)}}\quad(r\ge0),\qquad
 K_{\mathrm{nil}}=\Omega^{\,\bigcap_{c\ge1}\gamma_{c+1}(I)}.
                                                               \tag{META.DV5}
\]
Every finite extension of every $F_r$ and of
$K_{\mathrm{nil}}$ satisfies full closed-point-preserving
relative etale reconstruction for finite-type schemes
localized at universal homeomorphisms, with the
structural isomorphism retained.
The property also descends to any subfield over which
one of these reconstructed fields is algebraic.

In particular, for the Gauss $p$-valuation on
$D=\mathbb Q(t_1,\ldots,t_s)$, this gives positive
families with residue transcendence degree $s$ for
every $s\ge1$ and every prime $p$.
No assertion for all proper normal subgroups of
$G_M$ is included.
\end{claim}
\begin{proof}
We first establish the arithmetic inputs, and then
give the precise change in the analytic norm proof.
The tame/wild inertia theorem for an arbitrary DVR,
Stacks Tag~0BUA, Lemma~58.13.5, gives
\[
             1\longrightarrow P\longrightarrow I
               \longrightarrow \prod_{\ell\ne p}\mathbb Z_\ell(1)
               \longrightarrow1,                            \tag{META.DV6}
\]
with $P$ pro-$p$ \citep{AN_Stacks0BUA}.
The source does not require a perfect residue field.
All the prime-to-$p$ tame orders occur here by
adjoining roots of a uniformizer to the strict
henselization; units have roots of those orders by
Hensel's lemma.  The strict henselian residue field
is $k^{\mathrm{sep}}$, which need not be perfect.
Nothing below replaces it by an algebraic closure.

\medskip\noindent\textbf{The derived fields.}
Let $r\ge2$, put $K=F_r$ and $A=F_{r-1}$, and
let $E/K$ be finite.
If $a\in E^\times$ has zero rational Kummer class,
choose a finite Galois $D_1/D$ containing $a$ and
generators of $E/K$, and replace $E$ by $KD_1$.
Its completion $M_1$ is finite over $M$ and again
has finitely generated residue field.
Set
\[
 J=I\cap G_{D_1}=I_{M_1},\quad
 H=I^{(r)}\cap J,\quad \Gamma=J/H,\quad
 C=(I^{(r-1)}\cap J)/H,\quad V=\mathbb Q_p(1).
\]
The Kummer class on $J$ vanishes on $H$, which fixes
$V$, so inflates from a unique element of
$H^1(\Gamma,V)$.
For $g\in I^{(r-1)}$, Galoisness of $D_1/D$ makes
$g$ normalize $J$.
It acts trivially on $C$, on $\Gamma/C$ and on $V$.
Inflation--restriction gives a two-step filtration
of $H^1(\Gamma,V)$ on whose pieces $g$ is the identity:
\[
 0\longrightarrow H^1(\Gamma/C,V)
  \longrightarrow H^1(\Gamma,V)
  \longrightarrow H^1(C,V)^{\Gamma/C}.
\]
Thus $(g-1)^2=0$.
The particular Kummer vector has finite $g$-orbit,
because $a\in D_1$ has one.  If $g^n$ fixes it,
the equality $g^n-1=n(g-1)$ on that vector makes
$g$ fix it as well.
Therefore $\kappa_J(g(a)/a)=0$.
The inertia restriction injection and the finite
root-of-unity Kummer kernel of Claim~B.META.20 give
\[
       g(a)/a\in\mu(M_1)\quad(g\in I^{(r-1)}),\qquad
                         a^m\in A\text{ for some }m\ge1.       \tag{META.DV7}
\]
One exponent works for all $g$, since $g$ preserves
the one field $M_1$ and $\mu(M_1)$ is finite.
This proof includes every finite $E/K$.

By Claim~B.META.20 the derived inclusions are strict,
so choose $b\in K\setminus A$.
Both fields contain all roots of unity, since the
cyclotomic action of inertia is abelian.
If the classes of $b$ and $1-b$ both vanished,
(META.DV7) would give
$g(b)=\zeta_gb$ and $g(1-b)=\eta_g(1-b)$, with
$\zeta_g,\eta_g\in\mu_\infty\subset A$.
If some $\zeta_g\ne\eta_g$, adding expresses $b$
as an element of $A$; if all are equal, adding
makes them one, so again every $g$ fixes $b$.
This contradiction supplies a nonzero Kummer class
over $K$, surviving every finite extension by
rational restriction and corestriction.

For projection degree $d$ use the grids
\[
       \mathcal G_0=\{u+iv:0\le u,v\le4d\},
                         \qquad \mathcal G_1=b+\mathcal G_0.
\]
If a monic equal-degree rational function $R=P/Q$
has zero Kummer classes at all the grid values,
(META.DV7) makes $gR(a)/R(a)$ a root of unity at
every first-grid centre.
In a complex embedding of $\Omega$, the difference
of the squared absolute values of its monic
numerator and denominator has degree at most $4d$
separately in $u,v$.
Successive interpolation makes it zero.
Writing $z=u+iv,w=u-iv$ gives the product identity
$U(z)\overline U(w)=V(z)\overline V(w)$; monicity
then gives $gR=R$.
The translated grid similarly gives
$gR(gb+Z)=R(b+Z)$.
Choose $g$ moving $b$.
A rational function with the nonzero additive
period $gb-b$ is constant, by its finite pole set
or its leading polynomial finite difference.
Its value at infinity is one, so $R=1$.
This proves the required uniform two-grid test.

\medskip\noindent\textbf{The nilpotent union.}
Put $N=\bigcap_c\gamma_{c+1}(I)$ and $K=\Omega^N$.
Every finite inertia quotient is an extension of
a finite $p$-group by a cyclic prime-to-$p$ group.
In a nilpotent quotient all Sylow subgroups commute,
and all tame cyclic quotients survive.
Consequently
\[
       I/N=Q_{(p)}\times\prod_{\ell\ne p}\mathbb Z_\ell(1),
                                                               \tag{META.DV8}
\]
where $Q_{(p)}$ is pro-$p$.
Let $e=v_M(p)$ for the integer normalization.
Choose an odd prime $\ell\nmid ep(p-1)$, and put
$b=p^{1/\ell}$.
The strict henselization has value group
$e^{-1}\mathbb Z$ in the normalization $v(p)=1$;
thus adjoining $b$ gives a tame cyclic extension
of degree $\ell$.
Choose a central element in the pro-$\ell$ factor
of (META.DV8) moving $b$, and lift it to $\tau\in I$.
Its $p$-cyclotomic value is one: a pro-$\ell$ group
maps neither to a nontrivial pro-$p$ subgroup of
$\mathbb Z_p^\times$ nor to its torsion subgroup
of order dividing $p-1$.
Inertia fixes all other roots of unity, so $\tau$
fixes $\mu_\infty$ and $\tau b=\omega b\ne b$.

Here is the stable-kernel argument over this
complete DVR.
If $\kappa_E(a)=0$ for a finite $E/K$, then
compatible $p$-power roots of $a$ belong to $E$:
$G_E$ fixes $\mathbb Z_p(1)$, so its integral
$H^1$ is a torsion-free group of homomorphisms.
Choose finite Galois $D_1/D$ containing the data.
Let $F$ be the compositum of $D_1$, the relative
strict henselization, and $\mu_{p^\infty}$;
let $R/F$ be generated by all $p$-power roots of
the finitely many inertia conjugates of $a$.
Both fields are Galois over the strict henselization,
$R\subset KD_1$, and
$D_R=\operatorname{Gal}(R/F)$ is a torsion-free
closed subgroup of $\mathbb Z_p(1)^m$.
The image of $N$ in the Galois group of $R$
over the strict henselization is finite, since
$KD_1/K$ is finite.
Centrality of $\tau$ modulo $N$ puts
$[\tau,D_R]$ both in that finite image and in
$D_R$, hence makes it trivial.

For $J=I_{M_1}$ and $J_0=J\cap\ker\chi_p$,
the group $J_0$ surjects onto $D_R$.
The commutator equality, including the trivial
coefficient action of $\tau$, therefore gives
$\kappa_{J_0}(\tau a)=\kappa_{J_0}(a)$.
Restriction from $H^1(J,V)$ to $H^1(J_0,V)$ is
injective: $J/J_0$ is abelian with nontrivial
cyclotomic character, and comparing a cocycle
at $xy$ and $yx$ gives
$(\chi(x)-1)c(y)=(\chi(y)-1)c(x)$, so it is a
coboundary.
Claim~B.META.20 also gives inertia restriction
injectivity and the Kummer kernel in $M_1$.
We conclude
\[
             \kappa_E(a)=0\quad\Longrightarrow\quad
                                  \tau(a)/a\in\mu_\infty.      \tag{META.DV9}
\]
Choose the complex embedding with $b>0$.
Since
$|1-\omega b|^2-|1-b|^2=2b(1-\operatorname{Re}\omega)>0$,
the class of $1-b$ is nonzero and remains so over
every finite extension.
For $\mathcal G_0,b+\mathcal G_0$, the preceding
absolute-square interpolation applied to $\tau$
first gives $\tau R=R$ and then
$R(\tau b+Z)=R(b+Z)$.
The nonzero period gives $R=1$.
Thus both arithmetic inputs hold at every prime,
including $2$ and $3$.

\medskip\noindent\textbf{Local normalization and cyclic extensions.}
For either field $K$ just treated, start with an
arbitrary finite $E/K$, a smooth geometrically
integral affine source $X/E$, and a pinned relative
map to $\mathbb A^1_E$.
If its actual point map is nonconstant, each
actual fibre is a proper closed subset.
Choose a finite Noether projection of degree $d$
and the two grids above.
On a principal open $V$ of affine space avoiding
the branch locus and all the grid-fibre images,
the pulled-back and corestricted Kummer classes are
\[
          \xi_a=\beta_a+\sum_j\lambda_{aj}\kappa(h_j),
                      \qquad\lambda_{aj}\in\mathbb Q_p.        \tag{META.DV10}
\]
The geometric Picard group is zero, and the
geometric units modulo constants are freely
generated by the prime boundary equations $h_j$,
which proves this decomposition by Kummer theory
and inflation--restriction.
The actual evaluation is $\kappa(P_t(a))$ for
the monic degree-$d$ norm polynomial $P_t$;
the factor $(-1)^d$ from corestricting $Y-a$
is torsion and disappears in rational Kummer theory.

Choose one finite Galois $D_1/D$ containing
all coefficients, grid centres, a normalizing
parameter $t_*$ and its actual norm values,
all coordinates of the finitely many points in
one good projection fibre, and generators of $E/K$.
Let $M_1$ be its completion and $E_1=KD_1$.
The normality of the inertia kernel gives
$N_1=G_{E_1}\triangleleft G_{M_1}$ and
$\chi_p(N_1)=1$.
Evaluation at $t_*$ expresses
$\beta_a=\kappa(P_{t_*}(a))-\sum_j\lambda_{aj}
\kappa(h_j(t_*))$ as the restriction of a class
over $M_1$.
Thus (META.DV10), including its Tate coefficient
action, is $G_{M_1}$-invariant.
Comparison at $t$ and $\sigma t$ puts all grid
values of $P_{\sigma t}/\sigma P_t$ in the stable
kernel over a common finite extension.
The proved two-grid test yields
$P_{\sigma t}=\sigma P_t$ uniformly.

The nonzero Kummer class restricts nontrivially
to $N_1$.
Its integral image is $p^u\mathbb Z_p$; taking
an actual $p^u$-th root of its argument in $E_1$
makes that image all of $\mathbb Z_p$.
Adjoin this argument and $\mu_{p^s}$ to $M_1$
inside the local fixed field of $N_1$.
The resulting cyclic extension $L/M_1$ generated
by a $p^s$-th root has degree $q=p^s>d+1$ and
receives a surjection from $N_1$.
Normality persists under these finite enlargements.
Claim~B.META.20 supplies an $M_1$-basis of $L$
in $L\cap\Omega$, and
$E_1(L\cap\Omega)/E_1$ is finite:
the henselian $E_1$ contains $M_1\cap\Omega$
after the indicated algebraic enlargements.
All scalar algebraic parameter sets used below
are consequently dense.

Actual maps are valuation-continuous over every
finite extension of $E_1$.
Indeed CHW Proposition~3.6 and the representable
etale inverse-image comparison in \S4.13 and
Corollary~A.9 give etale-open continuity, and
JTWY Corollary~6.13 identifies that topology
with the valuation topology for these henselian
non-separably-closed fields
\citep{AN_CHW2024,AN_JTWY2024}.
The nonzero Kummer class established above
verifies non-separable-closedness at every finite
extension.  Affine charts handle nonseparated
finite-type schemes.

\medskip\noindent\textbf{The scalar-error step using deep units only.}
The chosen good etale projection fibre already
splits over $M_1$, before the cyclic $L/M_1$
was chosen.  Work near it with the $d$ analytic
inverse branches; no further base enlargement
is used to split this fibre.
For actual parameters in $L\cap\Omega$, local
norm equivariance puts the coefficients of $P_t$
in $L$; continuity gives continuity of those
coefficients on the dense algebraic parameter set.
At such a parameter put
\[
 \theta_t=\kappa_L(P_t(a))-\beta_a|_{G_L}
                  -\sum_j\lambda_{aj}\kappa_L(h_j(t)).
\]
It vanishes on $N_1\cap G_L$.
By inflation--restriction it inflates from
$G_L/(N_1\cap G_L)$.
Conjugation by $N_1$ is the identity on this
quotient and on $V$, since $N_1$ is normal in
$G_{M_1}$, $L/M_1$ is Galois, and $\chi_p(N_1)=1$.
Hence every $\theta_t$ is $N_1$-invariant.

Choose an actual centre $t_0$ and shrink the ball
so that, for every grid centre and every $h_j$,
the ratios $P_t(a)/P_{t_0}(a)$ and
$h_j(t)/h_j(t_0)$ lie in a deep unit group of $L$.
This uses only their actual continuity, not an
analytic extension of $P_t$.
Claim~B.META.20 now gives the exact identity
\[
 \theta_t-\theta_{t_0}
   =\iota_L\left(
      \log\frac{P_t(a)}{P_{t_0}(a)}
        -\sum_j\lambda_{aj}\log\frac{h_j(t)}{h_j(t_0)}
             \right).                                        \tag{META.DV11}
\]
Injectivity and Galois equivariance of $\iota_L$
make the scalar in parentheses fixed by $N_1$.
Since $N_1$ surjects onto $\operatorname{Gal}(L/M_1)$,
that scalar lies in $M_1$.
This is the exact replacement for a logarithm on
all of $H^1(G_L,V)$; residue-unit classes need not
belong to the deep-unit subspace.

Write $\tau$ for a generator of $\operatorname{Gal}(L/M_1)$
and $r_i(t,Z)=\tau^iP_t(Z)/P_t(Z)$ at actual points.
Subtract the $\tau^i$-transform of (META.DV11).
With the signs in the order transform minus original,
the result is
\[
 \log\frac{r_i(t,a)}{r_i(t_0,a)}
 =\sum_j\lambda_{aj}
   \left(
    \log\frac{h_j(\tau^it)}{h_j(\tau^it_0)}
       -\log\frac{h_j(t)}{h_j(t_0)}
   \right).                                                   \tag{META.DV12}
\]
After a further shrink, all finite linear combinations
are in the exponential convergence ball.
Exponentiating therefore recovers each ratio value
as $r_i(t_0,a)$ times an analytic function.
No global branch of logarithm on arbitrary residue
units has been chosen.
Differentiating gives exactly the matching-block
logarithmic equations of the norm proof in
Claim~B.META.14.

\medskip\noindent\textbf{Analytic recovery and the algebraic graph.}
We explain why that proof applies over the present
complete DVR without a local-compactness hypothesis.
Minimize the finite integer sum of conjugate gcd
degrees on the actual parameter ball.
At a minimizing point the relevant nonzero
subresultants remain nonzero by continuity;
the minimum then makes every gcd degree locally
constant on the dense actual set.
Recenter (META.DV11)--(META.DV12) at this
minimizing actual point.
For each reduced monic degree-$e_i$ ratio, its
values at $2e_i$ centres determine its numerator
and denominator analytically by (META.DV12).
The interpolation kernel is zero because
$\delta A B-A\delta B$ has degree at most $2e_i-1$
and has $2e_i$ roots; coprimality and monicity
then make both variations zero.

Use scalar coordinates for $\operatorname{Res}_{L/M_1}$,
with an algebraic $M_1$-basis and centre $t_0$.
After extending coefficients to $L$, the $q$
embedding blocks are permuted by the semilinear
cyclic action even if $t_0\notin M_1$.
The extended derivations of algebraic meromorphic
germs and the $4d+1$-centre rational identity
classify each nonconstant root by a unique block.
The rational identity has denominator degree at
most $4d$ and numerator degree at most $4d-1$,
so the centre count still proves it identically.
The divisor cocycle gives
$D_i=\tau^i B_0-B_0+C_i$.
A constant root is algebraic over $D$ after one
actual specialization, and its occurrence on an
open set contradicts proper closed actual fibres.
Thus $C_i=0$.
A negative coefficient of $B_0$ would give
$q-1>d$ distinct poles among the reduced ratios.
Their denominator least common multiple divides
the actual degree-$d$ polynomial at actual points,
which is impossible after a degree-preserving
specialization.

That specialization uses a nonempty analytic open,
not an algebraic open.
There are only finitely many nonzero minors and
subresultants to avoid.
A nonzero convergent power-series germ over a
complete nonarchimedean characteristic-zero field
cannot vanish on a scalar open ball: the
one-variable assertion follows by the first
nonzero term near an accumulation point, and
successive variables prove the multivariable
assertion.
Thus their complements are open and dense; a
finite intersection is nonempty and meets the
dense algebraic scalar parameter set.
Neither Baire's theorem nor local compactness
is used.

The effective type-zero divisor is already the
analytic monic reduced denominator $Q$ of $r_1$.
The exact derivatives obtained from (META.DV12)
give
$d\log Q(t,a)=\sum_j\lambda_{aj}\,dh_j/h_j$.
The constants of the meromorphic-germ differential
field are $L$, as follows term by term in a
power series.
Claim~B.135 therefore gives $Q\in L(t)[Z]$.
Its values at actual algebraic parameters are
algebraic over $D$, since it divides $P_t$ there;
finite rational interpolation gives
$Q\in\Omega(t)[Z]$.
Residues of its logarithmic derivatives make the
exponents integral.
Normalizing once more at an actual parameter
and using the two grids gives $P_t=Q_tC$ for a
single constant monic polynomial $C$.
Any root of a positive-degree $C$ would force
every projection fibre to meet one proper closed
actual fibre, contrary to finite-image dimension.
Hence $C=1$, and all projection norms have been
reconstructed.

The graph argument of Claim~B.META.14 is algebraic
over $\Omega$ and uses no additional local-field
property: one recovered norm bounds graph closure
finitely over a dense source open; a generic finite
linear projection separates its dominating
components and is birational on the generic
output fibres.
If their degrees over the source are $e_i$, the
distinct minimal output polynomials have degrees
$e_i d'$ and their product divides the recovered
degree-$d'$ norm.  Thus $\sum_i e_i\le1$, forcing
a unique component of degree one.
The finite coefficient conditions descend from
$\Omega(Y)$ to a nonempty open set over $\Omega$
by clearing denominators and expanding in $Y$.
Thus the actual map agrees densely with a rational
function, which descends to the original field
by equivariance.
The valuation continuity already proved excludes
poles using algebraic transverse curves, and
excludes exceptional point values by approximation
in an etale chart.
Normality of the smooth source makes the function
regular.
Geometric components descend after a finite
extension by the same uniqueness.
CHW Theorem~4.18 and Corollary~2.32 give the full
stated reconstruction property
\citep{AN_CHW2024}.

\medskip\noindent\textbf{Descent and strict additional scope.}
The proof has included every finite extension of
$F_r$ for $r\ge2$ and of $K_{\mathrm{nil}}$.
For an algebraic subfield, base change a pinned map
along the algebraic extension using the compatible
finite constant-object slices.
Every finitely presented etale object, morphism
and finite covering family descends to one finite
scalar extension, so these compatible slice functors
give the required inverse-image functor on the limit
site, preserving finite limits and covers.
Closed points remain closed since their residue
coordinates remain algebraic.
Reconstruct there and descend the unique resulting
scheme morphism by its equivariant actual point map.
This is the finite-site algebraic descent already
proved in Claim~B.META.16; it uses no finite-local-field
assumption.
Applying it to $F_0,F_1\subset F_2$ proves their
assertions, including every finite extension by
taking its finite compositum with $F_2$.

For $D=\mathbb Q(t_1,\ldots,t_s)$ the Gauss
valuation has group $\mathbb Z$ and residue
$\mathbb F_p(t_1,\ldots,t_s)$.
All fields in (META.DV5) are algebraic over $D$,
so have transcendence degree $s$ over $\mathbb Q$;
their residue fields are algebraic over that
residue field and hence have transcendence degree
$s$ over $\mathbb F_p$.
For $s>0$ they are not isomorphic to any henselian
subfield of $\overline{\mathbb Q}_p$ carrying its
induced $p$-adic valuation.
Indeed both valuations would be nontrivial of
rank one on the same non-separably-closed field.
If inequivalent they are independent, contradicting
the two-henselian-valuations theorem of F.\ K.\ Schmidt,
in the opened formulation Kuhlmann Theorem~9.44
\citep{AN_Kuhlmann9}; if equivalent their residue
transcendence degrees would agree, whereas every
residue field inside $\overline{\mathbb Q}_p$ is
algebraic over $\mathbb F_p$.
Thus this is a strictly additional field family.

The argument deliberately makes no assertion that
wild inertia for these residue fields is free,
or that all its nontrivial normal subgroups are
nonabelian.
Consequently it does not extend the all-proper-normal
conclusion of Claim~B.META.18 to the present bases.
That normal-subgroup boundary, other algebraic
subgroups, and more general nondiscrete valuations
remain separate questions.
No priority or external human verification is claimed.
\end{proof}
